% TOP_PROBLEM: {"id": 30007025, "title": "Existence of a finite ordinary Kähler fourfold program", "category_id": 5, "category": {"id": 5, "name": "algebraic_geometry", "display_name": "Algebraic Geometry", "description": "Geometric objects defined by polynomial equations.", "slug": "algebraic-geometry", "order_index": 5, "created_at": "2026-07-31T15:26:25.671Z"}, "status": "open"}
% FIRST_POSTED: 2026-10-06
% ATTEMPT_STATUS: solved
% OPENAI_CLAIM: {"schema_version":1,"source_repo":"https://github.com/openai/math","source_commit":"adc7f1241b42e322a6451854ab7e4b4c146bf78a","release_date":"2026-10-06","match":"full","resolution":"proved","independently_reviewed":false,"second_pass":"agreed","families":[{"family":"056","title":"Termination of projective and Kähler fourfold minimal model programs","subject":"Algebraic and complex geometry","manuscripts":[{"title":"Finite ordinary minimal model programs on compact Kähler fourfolds","pdf_path":"preprints/Finite-ordinary-minimal-model-programs-on-compact-Kahler-fourfolds-September-24-2026/paper.pdf","readme_path":"preprints/Finite-ordinary-minimal-model-programs-on-compact-Kahler-fourfolds-September-24-2026/README.md","date":"2026-09-24","theorem_ref":"thm:finite (sections/01-introduction.tex:97)","statement_sources":[{"path":"preprints/Finite-ordinary-minimal-model-programs-on-compact-Kahler-fourfolds-September-24-2026/build/sections/01-introduction.tex","line":97,"label":"thm:finite"}]}],"lean":null,"reasoning_trace":null}]}
\documentclass[11pt]{article}
\usepackage[margin=25mm]{geometry}
\usepackage{fontspec}
\setmainfont{Latin Modern Roman}
\usepackage{amsmath,amssymb,mathtools}
\usepackage{unicode-math}
\setmathfont{Latin Modern Math}
\usepackage{enumitem}
\usepackage{xurl,hyperref,bookmark}
\hypersetup{colorlinks=true,urlcolor=blue}
\setcounter{secnumdepth}{0}
\setlength{\parindent}{0pt}
\setlength{\parskip}{5pt}
\setlength{\emergencystretch}{3em}
\begin{document}
\section*{Existence of a finite ordinary Kähler fourfold program}

\subsection{OpenAI's claim}
For the stated compact Kähler klt fourfold pairs, at least one ordinary program is claimed to start on the supplied model, take finitely many steps, and finish nef or in a Mori fiber space. The cited manuscript argues for that result; its relation to this problem is specified below.

\subsection{Mathematical statement of the claimed result}
\paragraph{Mathematical claim 1.} For a normal globally Weil rational-factorial compact Kähler fourfold $X$ with rational effective $\Delta$ and klt pair $(X,\Delta)$, OpenAI claims existence of a finite ordinary negative-extremal-ray program starting on $X$, using projective divisorial contractions and flips. With $D=K_X+\Delta$, its endpoint satisfies
\[
D\text{ pseudo-effective}\Rightarrow D_n\text{ nef},\qquad
D\text{ not pseudo-effective}\Rightarrow X_n\to S\text{ a Mori fiber contraction}.
\]
All models remain normal globally Weil rational-factorial compact Kähler and klt. The first composite extracts no divisor and increases discrepancies weakly for all primes and strictly for contracted original primes. The Mori contraction has connected fibers, normal compact Kähler base of dimension below four, relative Picard rank one and relatively ample $-D_n$. Each contraction has relative Bott--Chern dimension one; strong rational-factoriality is preserved when assumed. Unlike the later manuscript, this main theorem states existence of a finite program, not termination of every choice.

\subsection{Outline of the argument}
Manuscript 1 (family 056, 2026-09-24): A single initial Kähler form is tracked as a positive current on the comparison models. Real-boundary terminal termination and power-free descent of actual lines preserve a fixed positive intersection gap, so rational perturbations transfer an auxiliary construction into the required ordinary program.

\subsection{Scope relative to this problem}
Manuscript 1 (family 056, 2026-09-24): The problem selects the mathematical statement summarized here, including its stated hypotheses and parameter range. The cited theorem gives a claimed proof of that assertion in full.

\subsection{Formal verification}
Family 056 has no Lean scope page in this snapshot; no formalization of these manuscript results is listed.

This site has not executed the Lean files or checked their dependencies. These descriptions report the scope of the linked declarations, rather than a new verification of the proofs.

\subsection{Status caveat}
OpenAI attributes these results to an internal model. Its release notice cautions that results lacking full formalization can contain mistakes. This site has not independently assessed the proofs; a solved label records the scope of the reported claim.

\subsection{Sources}
\begin{itemize}
\item OpenAI, \emph{Finite ordinary minimal model programs on compact Kähler fourfolds}, 2026-09-24; citation key \texttt{OAI:Finite-ordinary-minimal-model-programs-on-compact-Kahler-fourfolds-September-24-2026}. \href{https://github.com/openai/math/blob/adc7f1241b42e322a6451854ab7e4b4c146bf78a/preprints/Finite-ordinary-minimal-model-programs-on-compact-Kahler-fourfolds-September-24-2026/paper.pdf}{PDF (GitHub)}; \href{https://github.com/openai/math/raw/adc7f1241b42e322a6451854ab7e4b4c146bf78a/preprints/Finite-ordinary-minimal-model-programs-on-compact-Kahler-fourfolds-September-24-2026/paper.pdf}{PDF (direct)}; \href{https://github.com/openai/math/blob/adc7f1241b42e322a6451854ab7e4b4c146bf78a/preprints/Finite-ordinary-minimal-model-programs-on-compact-Kahler-fourfolds-September-24-2026/README.md}{Manuscript page and citation}; \href{https://github.com/openai/math/blob/main/preprints/Finite-ordinary-minimal-model-programs-on-compact-Kahler-fourfolds-September-24-2026/paper.pdf}{Latest version}; Statement source: \href{https://github.com/openai/math/blob/adc7f1241b42e322a6451854ab7e4b4c146bf78a/preprints/Finite-ordinary-minimal-model-programs-on-compact-Kahler-fourfolds-September-24-2026/build/sections/01-introduction.tex\#L97}{thm:finite}.
\item \href{https://github.com/openai/math/blob/adc7f1241b42e322a6451854ab7e4b4c146bf78a/CONTENTS.md}{OpenAI catalogue entry}.
\end{itemize}
\end{document}
